\documentclass[a4paper,10pt]{article}
\usepackage{tabularx}
\usepackage{geometry} % For page margins
\geometry{a4paper, margin=0.75in} % Set paper size and margins
\usepackage{titlesec} % For section title formatting
\usepackage{enumitem} % For compact itemized lists
\usepackage{hyperref} % For hyperlinks
\usepackage{etaremune}
\usepackage{libertine}
% Formatting
\setlength{\parskip}{0.5em} % Increase spacing between paragraphs and items
\titleformat{\section}{\large\bfseries}{1em}{0em}{}[\titlerule] % Section formatting
\titlespacing*{\section}{0pt}{1.5\baselineskip}{1.5\baselineskip} % Add extra space before and after sections
\titlespacing*{\subsection}{pt}{1.2\baselineskip}{0.8\baselineskip} % More space for subsections
% Define custom CV section command
\newcommand{\cvsection}[1]{\section*{\Large \textsc{#1}}\vspace{-1em}}
\newcommand{\cvsubsection}[1]{\noindent{\textsc{#1}}}

% Begin document
\begin{document}
% Header
\begin{center}
    {\LARGE \textbf{\sc{Han Li}}} \\[8pt] % Name
    Linguistics Department, Stony Brook University \\  
    100 Nicolls Road, Stony Brook, NY 11790 \\  
    \href{mailto:han.li.4@stonybrook.edu}{han.li.4@stonybrook.edu} |
    \href{http://lihan829.github.io}{lihan829.github.io}\\
    % \small {Updated \today}
\end{center}

% Education Section
\cvsection{Research Interests}
\begin{tabular}{@{}l p{0.8\linewidth}@{}}
\textbf{Theoretical} & Phonology · Phonetics · Tone and Prosody · Autosegmental Representations \\
\textbf{Computational} & Formal Language Theory · Learnability · Language Modeling \\
\textbf{Languages} & Hausa · Mandarin · Cantonese · Wu · Hokkien · Lanzhou · Uyghur 
\end{tabular}

\cvsection{Education}
\noindent
\textbf{Stony Brook University}, Stony Brook, NY \hfill 2021–Present \\  
\textit{Ph.D. in Linguistics} \\  
Affiliated with the \textbf{Institute for Advanced Computational Science} \\
Dissertation: \textit{Learning Tonal Patterns} \\
Committee: Jeffrey Heinz (main advisor), Ellen Broselow, Jiwon Yun, Adam Jardine (external)


\vspace{1em}\noindent
\textbf{Nanjing Aeronautics and Astronautics University}, Nanjing, China \hfill 2018–2021 \\  
\textit{M.A. in Linguistics} \\  
Thesis: \textit{Prosody of Mandarin Rhetorical Questions}  

\vspace{1em}\noindent
\textbf{China Agricultural University}, Beijing, China \hfill 2013–2017 \\  
\textit{B.A. in English Language and Literature} \\  
Thesis: \textit{English Interrogative Intonation of Uyghur EFL Learners} 
\\
\\
\noindent \textsc{Non-Degree Seeking}\\\\
\textbf{Institute for Advanced Computational Science;} Stony Brook, NY \hfill 2022-2024\\
Advanced Graduate Certificate in \textit{Data and Computational Science}
\cvsection{Publications}

\cvsubsection{Peer-reviewed Journals}

\begin{etaremune}[itemsep=0.8em] 
\item \textbf{Li, H.} \& Heinz, J. (Accepted) \textit{Tonotactic learning with autosegmental representations}. \textit{Phonology}.

\item Yi, L., \textbf{Li, H.}, Li, Y., \& Mu, J. (2024). An ongoing tonal-pattern change: 
Lanzhou dialect. \textit{Journal of Chinese Linguistics, 52}(2), 336–361.
\end{etaremune}

\cvsubsection{Peer-reviewed Conference Proceedings \& Book Chapters}
\begin{etaremune}[itemsep=0.8em]
\item \textbf{Li, H}. \& Heinz, J., (2026) What Matters in Tonotactic Learning. \textit{Society for Computation in Linguistics} 9(1).

\item \textbf{Li, H.} (2026). The Interaction between Tone and Intonation in Hausa Wh-questions. In Carstens, Vicki, Russell, Katherine R., Akingbade, Olawale, Morton, Deborah \& Diercks, Michael (eds.). 2026. \textit{Pamoja tena `Together again': African linguistics after COVID}. (Contemporary African Linguistics 14). Berlin: Language Science Press.


\end{etaremune}

\cvsubsection{Peer-reviewed Abstract-based Proceedings}
\begin{etaremune}[itemsep=0.8em]

\item \textbf{Li, H.} (2025). Learning tonotactic patterns over autosegmental representations. 
\textit{Proceedings of the Annual Meetings on Phonology 2025}

\item \textbf{Li, H.} (2021). Prosodic differences between rhetorical and information-seeking 
questions in Mandarin. \textit{ICU Working Papers in Linguistics, 15}, 45–52.

\end{etaremune}

% Presentations Section
\section*{\Large \textsc{Presentations}}
\vspace{-1em}
\cvsubsection{Conference and Workshop Presentations}
\begin{etaremune}[itemsep=0.8em]

    \item \textbf{Li, H.} \& Heinz, J. (2026). \textit{What Matters in Tonotactic Learning.} Talk at the Society for Computation in Linguistics, University of California,
    San Diego. 
    
    \item \textbf{Li, H.} \& Heinz, J. (2026) \textit{Comparing tonotactic learning over strings and autosegmental representations.} Old World Conference in Phonology (OCP 23), University of Cambridge, Cambridge, UK.

    
    \item \textbf{Li, H.} (2025). \textit{Understanding tones in language: patterns, representations, and insights.} Student Seminar, Institute for Advanced Computational Science, Stony Brook University.
    
    \item \textbf{Li, H.} (2025). \textit{Cross-fold Validation of BUFIA-AR}. Rutgers Subregular Phonology Workshop 2025, Rutgers University, New Brunswick, NJ.
    
    \item \textbf{Li, H.} (2024). \textit{Bottom-Up Factor Inference Algorithm over Autosegmental Representation (BUFIA-AR).} Rutgers Subregular Phonology Workshop, Rutgers University, New Brunswick, NJ.
    
    \item \textbf{Li, H.} (2024). \textit{Learning Tonal Patterns.} Lightning Talk, Institute for Advanced Computational Science, Stony Brook University.

    \item \textbf{Li, H.} (2022). \textit{The Q-morpheme in Hausa.} Paper presented at the 54th Annual Conference on African Linguistics (ACAL 54), University of Connecticut, Storrs, CT.

    \item \textbf{Li, H.} (2021). \textit{Prosodic differences between rhetorical and information-seeking questions in Mandarin.} Paper presented at the 5th Asian Junior Linguists Conference (AJL 5), International Christian University, Tokyo, Japan (online).

    \item \textbf{Li, H.} (2020). \textit{Syntactic distribution and prosodic features of rhetorical questions in Mandarin.} Paper presented at the 8th International Academic Conference for Graduates (IACGN 8), Nanjing, China.

    \item \textbf{Li, H.} (2019). \textit{Attenuation and boosting of complaints in Chinese computer-mediated communication.} Paper presented at the 7th International Academic Conference for Graduates (IACGN 7), Nanjing, China.

\end{etaremune}

\cvsubsection{Invited Talks}
\begin{etaremune}[itemsep=0.8em]
    \item \textbf{Li, H.} (2026). \textit{Tone as Structure: Morphology, Representation, and Learning}, Surrey Morphology Group, University of Surrey, London, UK.
    
    \item \textbf{Li, H.} (2025). \textit{Learning tonotactics over autosegmental representations.}, Phonology Group, Yale University, New Haven, CT.


\end{etaremune}

\cvsubsection{Poster Presentations}
\begin{etaremune}[itemsep=0.8em]

    \item \textbf{Li, H.} (2024). \textit{Learning tonotactic patterns over autosegmental representations.} Poster presented at the Annual Meeting on Phonology (AMP), Rutgers University, New Brunswick, NJ, October 2024.

    \item \textbf{Li, H.} (2024). \textit{Learning tonal patterns.} Poster presented at Institute for Advanced Computational Science Research Day, Stony Brook University, April 2024.

\end{etaremune}

% Honors & Awards Section

\cvsection{Honors \& Awards}
\begin{etaremune}[itemsep=0.3em]
\item Society for Computation in Linguistics (SCiL) Conference Scholarship (\$800) \hfill 2026
\item Stony Brook Foundation Board of Trustees Dissertation Completion Endowed Fellowship (\$15,000) \hfill 2026
\item Junior Researcher Award, Institute for Advanced Computational Science (\$34,000)  \hfill 2024 -2026
    \item Research Grant, Department of Linguistics, Stony Brook University (\$300)   \hfill 2021
    % \item Distinguished Graduate, College of Foreign Languages, Nanjing Aeronautics and Astronautics University \hfill 2021
    % \item Excellence Fellowship, College of Foreign Languages, Nanjing Aeronautics and Astronautics University \hfill 2018 - 2021
\end{etaremune}

% Teaching Experience Section
\cvsection{Employment}
\subsubsection*{\sc{Stony Brook University}}
\begin{itemize}[itemsep=0em, parsep=0.5em, topsep=0em]
\item \textbf{Instructor} - \textbf{Linguistics Department}  
\begin{itemize}[itemsep=0em, parsep=0.2em, topsep=0em]
    \item LIN 101   \quad Introduction to Linguistics (co-taught) \hfill Winter 2023
    \item LIN 230 \quad  Sociolinguistics (co-taught) \hfill Summer 2023; Summer 2024\\
\end{itemize}
\item \noindent\textbf{Teaching Assistant – Department of Linguistics} \\
{\footnotesize\textit{(Courses marked with * indicate responsibility for lab or recitation sessions.)}}

\begin{itemize}[itemsep=0em, parsep=0.3em, topsep=0em]
    \item LIN 637 \quad  Computational Linguistics \hfill Spring 2026
    \item LIN 330 \quad  Language Acquisition (Online, Asynchronous) \hfill Winter 2025 
    \item LIN 380 \quad  Anatomy and Physiology of Speech and Hearing \hfill Fall 2025
    \item LIN 350 \quad  Experimental Phonetics* \hfill Spring 2025
    \item LIN 320 \quad English Grammar \hfill Fall 2024
    \item LIN 110 \quad  Anatomy of Language (Online, Asynchronous) \hfill Spring 2024
    \item LIN 536 \quad Mathematical Methods \hfill Fall 2023
    \item LIN 201 \quad Phonetics* \hfill Fall 2022
    \item LIN 300 \quad Phonology* \hfill Spring 2022; Spring 2023
    \item LIN 431 \quad Uncommonly Taught Language (Hausa)* \hfill Fall 2021\\
\end{itemize}

\item \textbf{Research/Project Assistant – Department of Linguistics}\par
\begin{itemize}[itemsep=0em, parsep=0.3em, topsep=0em]
    \item Research Assistant to Dr.\ Thomas Graf, developing course materials for \textit{Math Methods} \hfill Summer 2022
    \item Research Assistant to Dr.\ John Bailyn at V-NYI  \hfill Spring 2026 (scheduled)\\
\end{itemize}

\item \textbf{Student Assistant – Student Accessibility \& Support Center} \hfill 2023–present
\begin{itemize}[itemsep=0em, parsep=0.3em, topsep=0em]
    \item Proctored accommodated exams and provided administrative support 
\end{itemize}
\end{itemize}

\cvsection{Service}  

\cvsubsection{Departmental / University}
\begin{itemize}[itemsep=0em, parsep=0.3em, topsep=0em]
    \item Student Coordinator, Open House, Stony Brook Linguistics PhD Program \hfill 2026
    \item Organizer, Stony Brook–Yale–NYU–CUNY Linguistics Conference (SYNC) \hfill 2024  
    \item Organizer, Linguistics Colloquium Series, Department of Linguistics \hfill 2023--2024
    \item Senator, Graduate Student Organization (GSO) \hfill 2022--2023
    \item Organizer, Linguistics Brown Bag Series, Department of Linguistics \hfill 2022\\
\end{itemize}

\cvsubsection{Professional Service}
\begin{itemize}[itemsep=0em, parsep=0.3em, topsep=0em]
    \item Reviewer, SYNC \hfill 2025
    \item Session Chair, Mathematics of Language (MoL 20) \hfill 2025
    \item Reviewer, Formal Approaches to South Asian Linguistics (FASAL 14) \hfill 2024
    \item Session Chair, Workshop on Model-Theoretic Representations in Phonology \hfill 2022\\
\end{itemize}

\cvsubsection{Outreach}
\begin{itemize}[itemsep=0em, parsep=0.3em, topsep=0em]
    \item Volunteer,  Longwood Women in STEM \hfill April 2026
    \item Instructor, North American Computational Linguistics Open Competition (NACLO) \hfill 2024--2025
    \item Tutor, Summer Youth Camp for Computational Linguistics (SYCCL) \hfill 2023--2025
\end{itemize}

% Skills Section
\cvsection{Skills}
\begin{itemize}[itemsep=0.3em]
    \item \textbf{Languages:} Mandarin (native), English (near-native), Japanese (intermediate)
    
    \item \textbf{Programming:} Python (Pandas, NumPy, scikit-learn, PyTorch), R (tidyverse), Haskell
    
    \item \textbf{Speech \& Experimental Tools:} Praat, ELAN, PsychoPy, E-Prime, Qualtrics
    
    \item \textbf{Computational \& Research Tools:} \LaTeX, Git, GitHub, Jupyter Notebooks
    
    \item \textbf{Methods:} Statistical analysis, experimental design, hypothesis testing, data visualization, machine learning (classification, regression, model evaluation)
\end{itemize}

\end{document}
